\chapter{Mission 4 and 5: Sound, Pictures, and Sharing}
\label{ch:kids3}

Your game works. Now make it \emph{feel} good: add sound, draw your own
pictures, and then put your game on a cartridge so that other people can play
it.

\section{Beep!}

The pink blocks in the \textbf{Audio} drawer make sound.

\kaudio{play beep freq 880 ms 80}

This block has two numbers.

\begin{itemize}
  \item \textbf{freq} is how \emph{high} the sound is. A big number is a high,
        squeaky sound. A small number is a low, growly sound.
  \item \textbf{ms} is how \emph{long} the sound lasts. 1000 is one second, so
        80 is very short.
\end{itemize}

\begin{prgmission}[Mission 4, step 1: happy beep, sad beep]
Open Lemon Catcher. Put a high beep inside the ``caught'' block and a low beep
inside the ``missed'' block:

\kinput{if rect lemon\_x lemon\_y 10 10 touches rect basket\_x 180 32 8}
\kstate[6mm]{change score by 1}
\kaudio[6mm]{play beep freq 880 ms 80}

\kinput{if rect lemon\_x lemon\_y 10 10 touches rect 0 196 320 4}
\kstate[6mm]{change misses by 1}
\kaudio[6mm]{play beep freq 220 ms 150}

Play with the sound on. Can you tell a catch from a miss with your eyes closed?
\end{prgmission}

\begin{prgwarn}
Put sound blocks \emph{inside an if}, so they play only when something
happens. A beep block sitting alone in \emph{every frame update} starts a new
beep every frame, and that is just noise.
\end{prgwarn}

\section{Notes}

The second pink block plays a musical note.

\kaudio{play MIDI note 60 on channel 0 for 250 ms}

Notes have numbers. Note 60 is the \emph{C} in the middle of a piano. Each step
up is the next key.

\begin{center}
\begin{tabular}{@{}lcccccccc@{}}
\toprule
Note name & C & D & E & F & G & A & B & C \\
Number & 60 & 62 & 64 & 65 & 67 & 69 & 71 & 72 \\
\bottomrule
\end{tabular}
\captionof{table}{The white keys from middle C up to the next C.}
\end{center}

\begin{prgmission}[Mission 4, step 2: the Music Pad]
Make a new project called \textbf{Music Pad}. Four buttons play four notes,
and a white bar shows which pad is playing.

\kgame{when game starts}
\kstate[6mm]{set marker\_x to 400}

\kgame{every frame update}
\kstate[6mm]{set marker\_x to 400}
\kinput[6mm]{if button LEFT pressed}
\kstate[12mm]{set marker\_x to 40}
\kaudio[12mm]{play MIDI note 60 on channel 0 for 120 ms}
\kinput[6mm]{if button UP pressed}
\kstate[12mm]{set marker\_x to 104}
\kaudio[12mm]{play MIDI note 64 on channel 0 for 120 ms}

Add two more: RIGHT plays note 67 with the marker at 168, and A plays note 72
with the marker at 232.

\kgame{every frame draw}
\kdraw[6mm]{clear screen BLACK}
\kdraw[6mm]{draw rectangle x 40 y 80 w 48 h 48 color BLUE}
\kdraw[6mm]{draw rectangle x 104 y 80 w 48 h 48 color GREEN}
\kdraw[6mm]{draw rectangle x 168 y 80 w 48 h 48 color RED}
\kdraw[6mm]{draw rectangle x 232 y 80 w 48 h 48 color MAGENTA}
\kdraw[6mm]{draw rectangle x marker\_x y 136 w 48 h 6 color WHITE}
\end{prgmission}

Why is marker\_x set to 400 at the start of every update? The screen is only
320 wide, so 400 is \emph{off the screen}. The bar is hidden there until a
button moves it under a pad. Hiding things off the screen is an old game
maker's trick.

\begin{prgtry}
Play \emph{C, E, G, C}: left, up, right, A. That is a chord called C major.
Change the note numbers to make your own four-note tune. Can a friend copy it
by ear?
\end{prgtry}

\section{Sprite Lab: draw your own pictures}

A small picture used in a game is called a \textbf{sprite}. Click
\textbf{Sprite Lab} at the top of the Kit.

\begin{prgmission}[Mission 4, step 3: a lemon with a face]
\begin{enumerate}
  \item Click \kbtn{New Sprite}. Name it \textbf{lemon}. Keep 16 by 16.
  \item Pick yellow and paint a round lemon, one pixel at a time.
  \item Pick black and give it two eyes and a smile.
  \item Made a mistake? Use the eraser. Erased pixels are see-through.
  \item Click \kbtn{Save}.
\end{enumerate}
\end{prgmission}

A 16 by 16 sprite has $16 \times 16 = 256$ pixels. Every pixel is one colour.
Up close you see squares. From far away you see a lemon. All pictures on all
screens work like this.

\section{What is under the blocks?}

Here is a secret. The computer does not really understand blocks. It
understands a language of words and symbols. One such language is called
\textbf{C}. Grown-up programmers type it by hand.

Click \kbtn{Generate C}. On the right you can see your game written in C.

\begin{center}
\begin{tabularx}{\linewidth}{@{}XX@{}}
\toprule
Your block & The same thing in C \\
\midrule
set score to 0 & \code{score = 0;} \\
change basket\_x by 4 & \code{basket\_x += 4;} \\
if button LEFT pressed & \code{if (input \& PRG32\_BTN\_LEFT) \{ \}} \\
clear screen BLACK & \code{prg32\_gfx\_clear(PRG32\_COLOR\_BLACK);} \\
\bottomrule
\end{tabularx}
\captionof{table}{Each block becomes one line of C.}
\end{center}

\begin{prgtry}
Find the line in the C that came from your \textbf{keep basket\_x between}
block. How many \textbf{if}s did one block turn into?
\end{prgtry}

You do not have to understand all of it. Just notice: your blocks \emph{are} a
real program. One day you may write the words yourself. The older students who
use this book do exactly that.

\section{Mission 5: your game on a cartridge}

Long ago, video games came on \textbf{cartridges}: plastic boxes you pushed
into the machine. \prg games come on cartridges too, but ours are files. A \prg
cartridge file ends in \fname{.prg32}.

\begin{prgmission}[Mission 5, step 1: make the cartridge]
Open your best game. Click \kbtn{Compile Cartridge}.

If your teacher has set up the cartridge tools, green download buttons appear.
You have made a real cartridge!
\end{prgmission}

\begin{center}
\begin{tikzpicture}[font=\sffamily\small,node distance=7mm,
  s/.style={rounded corners=3pt,draw=prgsteel,fill=prgpanel,minimum height=11mm,minimum width=23mm,align=center},
  ar/.style={-{Stealth[length=2.2mm]},draw=prgink!70,thick}]
  \node[s,fill=prgamber!20] (b) {your\\blocks};
  \node[s,right=of b] (c) {C\\program};
  \node[s,right=of c] (k) {cartridge\\\fname{.prg32}};
  \node[s,right=of k,fill=prggreen!15] (m) {a real\\\prg};
  \draw[ar] (b) -- (c); \draw[ar] (c) -- (k); \draw[ar] (k) -- (m);
\end{tikzpicture}
\captionof{figure}{From blocks to a real machine.}
\end{center}

A cartridge can be played in three places:

\begin{itemize}
  \item on a real \textbf{\prg console}, a small computer with a screen and a
        joystick;
  \item in \textbf{PRG32-QT}, an app for computers, tablets, phones, and TVs;
  \item by anyone who finds it in the \textbf{Cartridge Store}.
\end{itemize}

\section{The Cartridge Store}

The Cartridge Store is a web site full of \prg games, like a library. Anyone
can look and play. To \emph{add} a game, you send it in, and a person called an
\textbf{editor} checks it first. In your class, the editor is probably your
teacher.

\begin{prgmission}[Mission 5, step 2: share, with your teacher]
\begin{enumerate}
  \item Give your game a good title and a short sentence that says what it is.
  \item Ask your teacher to open the \kbtn{Publish} page with you.
  \item Click \kbtn{Publish} next to your game.
  \item Wait. Your teacher checks the game and says yes.
  \item Find your game in the Store. Show a friend!
\end{enumerate}
\end{prgmission}

\begin{prgconcept}
Sharing has rules. Share only what \emph{you} made. Never put your full name,
your address, or your photo in a game. Be kind in titles and words: real people
will read them.
\end{prgconcept}

\section{What you have learned}

Look how far you have come. You know the big ideas that every programmer
uses, in every language:

\begin{center}
\begin{tabularx}{\linewidth}{@{}lX@{}}
\toprule
Big idea & Where you used it \\
\midrule
sequence & blocks run from top to bottom \\
loop & update and draw repeat every frame \\
variable & player\_x, score, misses \\
condition & if button pressed, if touching \\
event & a beep when a lemon is caught \\
coordinates & x across, y down \\
randomness & a new lemon somewhere new \\
debugging & Step, and watching the numbers \\
\bottomrule
\end{tabularx}
\end{center}

\begin{prgcheck}
You finished the Young Makers track if you can:
\begin{itemize}
  \item add a sound that plays only when something happens;
  \item draw and save a sprite;
  \item find one of your blocks in the generated C;
  \item explain what a cartridge is;
  \item say two rules for sharing safely.
\end{itemize}
\end{prgcheck}

\begin{prggrownup}
\kbtn{Compile Cartridge} produces real \fname{.prg32} files only when the \prg
toolchain is installed where the Kit runs; otherwise it produces a source-only
bundle (Chapter~\ref{ch:kitteach}). Publication always goes through an editor's
review (Chapter~\ref{ch:store}); pupils should never hold publication tokens.
Sprite Lab converts drawings to C arrays; using a sprite inside a Blocks game
requires the Advanced~C editor, which is a good extension for confident pupils
aged eleven and over.
\end{prggrownup}
